% Options for packages loaded elsewhere
\PassOptionsToPackage{unicode}{hyperref}
\PassOptionsToPackage{hyphens}{url}
\documentclass[
]{article}
\usepackage{xcolor}
\usepackage[margin=1in]{geometry}
\usepackage{amsmath,amssymb}
\setcounter{secnumdepth}{-\maxdimen} % remove section numbering
\usepackage{iftex}
\ifPDFTeX
  \usepackage[T1]{fontenc}
  \usepackage[utf8]{inputenc}
  \usepackage{textcomp} % provide euro and other symbols
\else % if luatex or xetex
  \usepackage{unicode-math} % this also loads fontspec
  \defaultfontfeatures{Scale=MatchLowercase}
  \defaultfontfeatures[\rmfamily]{Ligatures=TeX,Scale=1}
\fi
\usepackage{lmodern}
\ifPDFTeX\else
  % xetex/luatex font selection
\fi
% Use upquote if available, for straight quotes in verbatim environments
\IfFileExists{upquote.sty}{\usepackage{upquote}}{}
\IfFileExists{microtype.sty}{% use microtype if available
  \usepackage[]{microtype}
  \UseMicrotypeSet[protrusion]{basicmath} % disable protrusion for tt fonts
}{}
\makeatletter
\@ifundefined{KOMAClassName}{% if non-KOMA class
  \IfFileExists{parskip.sty}{%
    \usepackage{parskip}
  }{% else
    \setlength{\parindent}{0pt}
    \setlength{\parskip}{6pt plus 2pt minus 1pt}}
}{% if KOMA class
  \KOMAoptions{parskip=half}}
\makeatother
\setlength{\emergencystretch}{3em} % prevent overfull lines
\providecommand{\tightlist}{%
  \setlength{\itemsep}{0pt}\setlength{\parskip}{0pt}}
\usepackage{bookmark}
\IfFileExists{xurl.sty}{\usepackage{xurl}}{} % add URL line breaks if available
\urlstyle{same}
\hypersetup{
  hidelinks,
  pdfcreator={LaTeX via pandoc}}

\author{}
\date{}

\begin{document}

\section{Mathematical \& Methodological Architecture of the PhD
Dissertation: Rigorous Multi-Metric Benchmarking of Model-Agnostic
Explainable
AI}\label{mathematical-methodological-architecture-of-the-phd-dissertation-rigorous-multi-metric-benchmarking-of-model-agnostic-explainable-ai}

\emph{Advanced PhD Level (Researchers, Doctoral Committee, \& Peer
Reviewers)}

\begin{center}\rule{0.5\linewidth}{0.5pt}\end{center}

\subsection{1. Formal Problem Formulation \& Theoretical
Foundations}\label{formal-problem-formulation-theoretical-foundations}

Let \(f: \mathcal{X} \to \mathcal{Y}\) represent an opaque, non-linear
black-box predictor trained over an input feature space
\(\mathcal{X} \subseteq \mathbb{R}^d\) to map to a target space
\(\mathcal{Y}\). In local post-hoc Explainable AI (XAI), given a target
instance \(\mathbf{x} \in \mathcal{X}\), the goal is to construct an
interpretable model \(g \in \mathcal{G}\) operating over a simplified
binary input vector \(\mathbf{x}' \in \{0, 1\}^M\) that locally
approximates the decision boundary of \(f\) around \(\mathbf{x}\).

\subsubsection{1.1 Mathematical Formulation of Post-Hoc Explainer
Paradigms}\label{mathematical-formulation-of-post-hoc-explainer-paradigms}

\paragraph{A. Local Interpretable Model-agnostic Explanations
(LIME)}\label{a.-local-interpretable-model-agnostic-explanations-lime}

LIME formulates local explanation as an additive feature attribution
optimization problem:
\[\arg\min_{g \in \mathcal{G}} \mathcal{L}(f, g, \pi_{\mathbf{x}}) + \Omega(g)\]
where \(\mathcal{L}\) measures the local unfaithfulness of \(g\) in
approximating \(f\) over a perturbed sample space \(\mathcal{Z}\),
weighted by an exponential proximity kernel:
\[\pi_{\mathbf{x}}(\mathbf{z}) = \exp\left( -\frac{D(\mathbf{x}, \mathbf{z})^2}{\sigma^2} \right)\]
and \(\Omega(g)\) penalizes model complexity (e.g., constraining the
\(L_0\) norm \(\|g\|_0 \le K\)).

\paragraph{B. Shapley Additive exPlanations (SHAP \&
KernelSHAP)}\label{b.-shapley-additive-explanations-shap-kernelshap}

SHAP calculates the unique additive feature attribution vector
\(\boldsymbol{\phi} = (\phi_1, \dots, \phi_M) \in \mathbb{R}^M\)
satisfying four fundamental game-theoretic axioms: \textbf{Efficiency}
(\(\sum \phi_i = f(\mathbf{x}) - \mathbb{E}[f(\mathbf{x})]\)),
\textbf{Symmetry}, \textbf{Dummy/Nullity}, and \textbf{Additivity}.

The classical Shapley value for feature \(i\) given characteristic
function
\(v(S) = \mathbb{E}_{\mathbf{z} \sim \mathcal{D}}[f(\mathbf{z}) \mid \mathbf{z}_S = \mathbf{x}_S]\)
is defined by:
\[\phi_i(v) = \sum_{S \subseteq N \setminus \{i\}} \frac{|S|!(|N|-|S|-1)!}{|N|!} \Big[ v(S \cup \{i\}) - v(S) \Big]\]

KernelSHAP estimates \(\boldsymbol{\phi}\) by solving a weighted linear
regression over simplified coalition vectors
\(\mathbf{z}' \in \{0,1\}^M\) with Shapley kernel weights:
\[\pi_{\mathbf{x}}(\mathbf{z}') = \frac{M - 1}{\binom{M}{|\mathbf{z}'|} |\mathbf{z}'| (M - |\mathbf{z}'|)}\]

\paragraph{C. Anchors (High-Precision Rule
Explanations)}\label{c.-anchors-high-precision-rule-explanations}

An Anchor is a rule predicate \(A \subseteq \text{predicates}\) such
that \(A(\mathbf{x}) = 1\), maximizing coverage
\(\text{cov}(A) = \mathbb{P}_{\mathbf{z} \sim \mathcal{D}}[A(\mathbf{z}) = 1]\)
subject to a high-precision probabilistic constraint:
\[\mathbb{P}_{\mathbf{z} \sim \mathcal{D}(\cdot \mid A(\mathbf{z})=1)} \Big[ f(\mathbf{z}) = f(\mathbf{x}) \Big] \ge \tau \quad \text{with confidence } 1 - \delta\]
where \(\tau \in (0, 1]\) (typically \(\tau = 0.95\)).

\paragraph{D. Diverse Counterfactual Explanations
(DiCE)}\label{d.-diverse-counterfactual-explanations-dice}

DiCE generates a set of \(K\) diverse counterfactual instances
\(\mathbf{C} = \{\mathbf{c}_1, \dots, \mathbf{c}_K\}\) solving the
multi-objective loss:
\[\arg\min_{\mathbf{c}_1, \dots, \mathbf{c}_K} \frac{1}{K}\sum_{k=1}^K \text{loss}(f(\mathbf{c}_k), y^*) + \frac{\alpha}{K}\sum_{k=1}^K \text{dist}(\mathbf{x}, \mathbf{c}_k) - \beta \cdot \text{DPP\_diversity}(\mathbf{c}_1, \dots, \mathbf{c}_K)\]
where \(y^* \neq f(\mathbf{x})\) is the desired counterfactual target
class, \(\text{dist}(\cdot, \cdot)\) is a normalized distance metric,
and \(\text{DPP\_diversity}\) maximizes the determinant of the diversity
kernel matrix using a Determinantal Point Process.

\begin{center}\rule{0.5\linewidth}{0.5pt}\end{center}

\subsection{2. The FOM-7 Operational Protocol: 7-Gate Governance
Architecture}\label{the-fom-7-operational-protocol-7-gate-governance-architecture}

To eliminate experimental drift, un-reproducible seed variation, and
reporting bias in XAI evaluation, this dissertation introduces the
\textbf{FOM-7 Protocol} (\emph{Framework for Operational Metrics in 7
Gates}).

\begin{verbatim}
[ Gate 1: Freezing ] ──> Pre-registered design & hypothesis matrix (300 cells)
         │
         ▼
[ Gate 2: Execution ] ──> Deterministic execution under manifest isolation
         │
         ▼
[ Gate 3: Audit ]     ──> Artifact integrity check (|E*| = 275 qualified cells)
         │
         ▼
[ Gate 4: Harmonize ] ──> Block-level aggregation over K=15 (model x sample-size) blocks
         │
         ▼
[ Gate 5: Export ]    ──> Non-parametric inferential testing (Friedman + Nemenyi + Holm)
         │
         ▼
[ Gate 6: Profile ]   ──> Variance decomposition (CV = sigma/mu) separating seeds from methods
         │
         ▼
[ Gate 7: Report ]    ──> Provenance-backed claim registry (pub/claim_registry.toml)
\end{verbatim}

\subsubsection{2.1 Formal Specification of the 7
Gates}\label{formal-specification-of-the-7-gates}

\begin{enumerate}
\def\labelenumi{\arabic{enumi}.}
\tightlist
\item
  \textbf{Gate 1 (Freezing):} Pre-registration of experimental space
  \(\mathcal{E} = \mathcal{M} \times \mathcal{X}_{exp} \times \mathcal{S} \times \mathcal{N}\),
  where \(|\mathcal{M}|=5\) model families
  (\(\text{LogReg}, \text{RF}, \text{XGB}, \text{SVM}, \text{MLP}\)),
  \(|\mathcal{X}_{exp}|=4\) explainers
  (\(\text{LIME}, \text{SHAP}, \text{Anchors}, \text{DiCE}\)),
  \(|\mathcal{S}|=5\) random seeds, and \(|\mathcal{N}|=3\) sample sizes
  (\(N \in \{100, 500, 1000\}\)), producing 300 planned experimental
  cells.
\item
  \textbf{Gate 2 (Execution):} Containerized execution via deterministic
  manifest files.
\item
  \textbf{Gate 3 (Audit):} Rigorous filtering of cell artifacts:
  \[\mathcal{E}^* = \{ e \in \mathcal{E} \mid \text{status}(e) \in \{\text{ok\_instance}, \text{ok\_aggregated}\} \}\]
  In our benchmark, \(|\mathcal{E}^*| = 275\) qualified cells
  (\(91.67\%\) yield); non-qualifying cells (primarily constraint
  timeouts in DiCE/Anchors) are formally audited and logged rather than
  imputed.
\item
  \textbf{Gate 4 (Harmonization):} Block-level metric aggregation across
  \(K=15\) model-size blocks
  (\(\text{block}_{m, n} = \text{model}_m \times \text{size}_n\)).
\item
  \textbf{Gate 5 (Exportation):} Inferential statistical evaluation
  using non-parametric Friedman tests:
  \[Q = \frac{12 K}{k(k+1)} \left[ \sum_{j=1}^k \bar{R}_j^2 - \frac{k(k+1)^2}{4} \right] \sim \chi^2_{k-1}\]
  Followed by post-hoc Nemenyi tests with Critical Distance (\(CD\)):
  \[CD = q_{\alpha, k} \sqrt{\frac{k(k+1)}{6K}}\] with Holm-Bonferroni
  step-down multiplicity correction.
\item
  \textbf{Gate 6 (Profiling):} Variance decomposition isolating
  seed-induced stochasticity (\(\sigma^2_{\text{seed}}\)) from true
  explainer main effects (\(\sigma^2_{\text{method}}\)) via the
  Coefficient of Variation (\(CV = \frac{\sigma}{\mu}\)).
\item
  \textbf{Gate 7 (Reporting):} Automated verification mapping published
  numbers back to raw execution artifacts via
  \texttt{pub/claim\_registry.toml}.
\end{enumerate}

\begin{center}\rule{0.5\linewidth}{0.5pt}\end{center}

\subsection{3. The PhD Publication Matrix: Scientific Arc \&
Contributions}\label{the-phd-publication-matrix-scientific-arc-contributions}

\begin{verbatim}
                          ┌─────────────────────────────────────────┐
                          │   PhD Thesis: Rigorous XAI Evaluation   │
                          └────────────────────┬────────────────────┘
                                               │
         ┌─────────────────────────────────────┼─────────────────────────────────────┐
         │                                     │                                     │
         ▼                                     ▼                                     ▼
[ Theoretical & Methodological ]      [ Empirical Benchmark ]               [ Governance & Human ]
  ├── Paper A (RIMI 2026): FOM-7        ├── Paper B (CLEIej): LIME vs SHAP    ├── Paper E (CyS): Pareto Frontier
  └── Paper C (Tec. Marcha): Scaling    └── Paper D (Tec. Marcha): Adult      └── Paper F (JCSI): Human/LLM Judge
\end{verbatim}

\begin{center}\rule{0.5\linewidth}{0.5pt}\end{center}

\subsubsection{Paper A: Theoretical \& Methodological Framework (FOM-7
Protocol)}\label{paper-a-theoretical-methodological-framework-fom-7-protocol}

\begin{itemize}
\tightlist
\item
  \textbf{Venue \& Status:} Published in \emph{Revista de Investigación
  Multidisciplinaria Iberoamericana (RIMI)} (2026, DOI:
  \texttt{10.69850/rimi.vi3.307}).
\item
  \textbf{Core Contribution:} Establishes the formal mathematical
  foundation of multi-metric evaluation
  (\(\text{Fidelity}, \text{Stability}, \text{Sparsity}, \text{Cost}, \text{Faithfulness Gap}\))
  and specifies the 7-gate FOM-7 protocol.
\item
  \textbf{Key Finding:} Proves that single-metric evaluations (e.g.,
  fidelity alone) create false rankings because methods optimize
  different objective trade-offs.
\end{itemize}

\begin{center}\rule{0.5\linewidth}{0.5pt}\end{center}

\subsubsection{Paper B: Scaled Empirical Comparison --- LIME vs
SHAP}\label{paper-b-scaled-empirical-comparison-lime-vs-shap}

\begin{itemize}
\tightlist
\item
  \textbf{Venue \& Status:} Submitted to \emph{CLEI Electronic Journal
  (CLEIej)} (2026-10-04, Subm. \#1196).
\item
  \textbf{Core Hypotheses Tested:}

  \begin{itemize}
  \tightlist
  \item
    \(H_1\):
    \(\text{Stability}(\text{SHAP}) > \text{Stability}(\text{LIME})\)
    (\(p < 0.001\), confirmed).
  \item
    \(H_2\):
    \(\text{Sparsity}(\text{LIME}) < \text{Sparsity}(\text{SHAP})\)
    (confirmed).
  \item
    \(H_3\): \(\text{Cost}(\text{LIME}) \ll \text{Cost}(\text{SHAP})\)
    (\(p < 0.001\), confirmed).
  \end{itemize}
\item
  \textbf{Key Finding:} Proves that KernelSHAP's Shapley kernel
  optimization guarantees numerical stability under input perturbation,
  whereas LIME's monte-carlo perturbation sampling introduces
  significant seed variance (\(\sigma^2_{\text{seed}}\)), leading to
  instability across runs.
\end{itemize}

\begin{center}\rule{0.5\linewidth}{0.5pt}\end{center}

\subsubsection{Paper C: Scalability, Sample Size Dynamics, \&
Robustness}\label{paper-c-scalability-sample-size-dynamics-robustness}

\begin{itemize}
\tightlist
\item
  \textbf{Venue \& Status:} Ready to submit to \emph{Tecnología en
  Marcha} (AI Special Issue, Zenodo v0.12.0 DOI:
  \texttt{10.5281/zenodo.23165763}).
\item
  \textbf{Core Contribution:} Analyzes the functional scaling behavior
  of explainers as background sample size \(N\) increases from 100 to
  1,000 instances.
\item
  \textbf{Key Finding:} Demonstrates that KernelSHAP's computational
  complexity scales as \(\mathcal{O}(N \cdot M \cdot 2^M)\) without
  background sub-sampling, creating a steep execution cost barrier,
  whereas tree-specific algorithms (TreeSHAP) exploit structural paths
  to reduce complexity to \(\mathcal{O}(N \cdot L \cdot D^2)\).
\end{itemize}

\begin{center}\rule{0.5\linewidth}{0.5pt}\end{center}

\subsubsection{Paper D: Empirical Validation on Tabular Benchmarks (UCI
Adult)}\label{paper-d-empirical-validation-on-tabular-benchmarks-uci-adult}

\begin{itemize}
\tightlist
\item
  \textbf{Venue \& Status:} Ready to submit to \emph{Tecnología en
  Marcha} (AI Special Issue).
\item
  \textbf{Core Contribution:} First comprehensive multi-model,
  multi-explainer benchmark execution over the \emph{UCI Adult Income}
  tabular dataset across 300 experimental cells.
\item
  \textbf{Key Finding:} Provides empirical mean ranks across 15 complete
  model blocks:

  \begin{itemize}
  \tightlist
  \item
    \textbf{Fidelity:} SHAP (\(\bar{R}=0.8081\)) significantly
    outperforms LIME (\(\bar{R}=0.5602\), \(p < 0.05\)).
  \item
    \textbf{Structural vs Attributional Methods:} Anchors and DiCE
    cannot be evaluated on linear feature attribution scales; they
    require rule-coverage and counterfactual distance metrics
    respectively.
  \end{itemize}
\end{itemize}

\begin{center}\rule{0.5\linewidth}{0.5pt}\end{center}

\subsubsection{Paper E: Pareto Optimization Frontiers \& XAI
Governance}\label{paper-e-pareto-optimization-frontiers-xai-governance}

\begin{itemize}
\tightlist
\item
  \textbf{Venue \& Status:} Submitted to \emph{Computación y Sistemas
  (CIC-IPN)} (2026-10-04, Subm. \#6783).
\item
  \textbf{Core Contribution:} Formulates XAI selection as a
  Multi-Objective Optimization Problem (MOOP):
  \[\max_{g \in \mathcal{G}} \Big( \text{Fidelity}(g), \text{Stability}(g), -\text{Cost}(g), -\text{Sparsity}(g) \Big)\]
\item
  \textbf{Key Finding:} Identifies the non-dominated Pareto frontier:

  \begin{itemize}
  \tightlist
  \item
    \textbf{Healthcare / High-Stakes Audit:} SHAP dominates (maximal
    stability and fidelity).
  \item
    \textbf{Low-Latency Production API:} LIME / TreeSHAP dominate
    (minimal cost, low sparsity).
  \item
    Provides a formal governance decision framework for deploying XAI in
    regulated industries.
  \end{itemize}
\end{itemize}

\begin{center}\rule{0.5\linewidth}{0.5pt}\end{center}

\subsubsection{Paper F: Human \& LLM-as-a-Judge Semantic
Alignment}\label{paper-f-human-llm-as-a-judge-semantic-alignment}

\begin{itemize}
\tightlist
\item
  \textbf{Venue \& Status:} Early development (\emph{Journal of Computer
  Sciences Institute}).
\item
  \textbf{Core Contribution:} Investigates the semantic alignment
  between quantitative mathematical metrics
  (\(\text{Fidelity}, \text{Stability}\)) and human/LLM cognitive
  evaluations.
\item
  \textbf{Key Finding:} Formulates a pre-analysis plan testing whether
  high-fidelity mathematical attribution vectors correspond to superior
  human task performance and LLM explanation preference.
\end{itemize}

\begin{center}\rule{0.5\linewidth}{0.5pt}\end{center}

\subsection{4. Synthesis \& Doctoral Defense
Narrative}\label{synthesis-doctoral-defense-narrative}

This dissertation unifies these six papers into a cohesive defense
thesis:

\[\text{Opaque Model } f(\mathbf{x}) \xrightarrow[\text{Gate 1--3}]{\text{FOM-7 Protocol}} \text{Audited Explanations } g(\mathbf{x}) \xrightarrow[\text{Gate 4--6}]{\text{Pareto Evaluation}} \text{Defensible XAI Governance}\]

\begin{enumerate}
\def\labelenumi{\arabic{enumi}.}
\tightlist
\item
  \textbf{The Core Gap:} The unvalidated proliferation of post-hoc
  explainers created a crisis of evaluation in XAI.
\item
  \textbf{The Methodological Solution (Paper A):} The 7-gate FOM-7
  protocol provides mathematical rigor, statistical multiplicity
  control, and artifact auditability.
\item
  \textbf{Empirical Benchmarking (Papers B, C, D):} Extensive empirical
  evidence on 300 cells establishes exact performance trade-offs between
  LIME, SHAP, Anchors, and DiCE across 5 model families.
\item
  \textbf{Decision Governance \& Human Alignment (Papers E, F):} Pareto
  frontier analysis and semantic alignment bridge the gap between
  mathematical evaluation and real-world deployment.
\end{enumerate}

\end{document}
